\documentclass[11pt]{article}
\usepackage{mathblatt}
% Präambel des Setzers (werkzeuge/setzer-praeambel.tex, wortgleich); dazu \gteil (Muster T2U)
\makeatletter
\setlength{\mbpfbe}{0mm}% keine Punkte (pruefung.md § 5: nur beim Prüfstein)
\newlength{\gEin}\setlength{\gEin}{0pt}
\let\gSchrift\relax
\renewcommand{\pfkreuzzeile}[1]{\par\noindent\hangindent3.2em\hangafter1\hspace*{1.6em}$\square$\ #1\par}
\newcommand{\gkreuz}[1]{$\square$\ #1\hspace{2em}}
\newcommand{\gkreuzzeile}[2]{\par\noindent\hangindent2.2em\hangafter1\mbox{#1\,$\square$}\ #2\par}% Buchstabe, Option: eine je Zeile (neben Skizze)
% links, was man ansieht, rechts, was man tut; Spaltengrenze überall an derselben Stelle (gemeinsam.md „Satz“)
% v0.17: Fußeintrag „Nr: Wert“ – kurz in einer Box (nie getrennt), lang umbrechbar mit der Nummer am Anfang
% (vorher \mbox: lange Lösungen liefen über den Rand, AS40 138 pt); robust, weil \fusshilfe in die .aux schreibt
\NewDocumentCommand{\gfuss}{m m}{\sbox0{#1:~#2}\ifdim\wd0<0.45\textwidth\mbox{#1:~#2}\else\mbox{#1:}\nobreak\ #2\fi}
\newsavebox{\gbox}\newlength{\gLinks}
% v0.10 (Lehrer 13.10.): keine starre Halbierung – die linke Spalte ist so breit wie die Figur (+ 6 mm), die
% Antwort steht direkt daneben; ist die Figur breiter als 0,62 Zeile, steht die Antwort darunter
\newcommand{\gzwei}[2]{\par\smallskip\noindent\sbox{\gbox}{#1}%
  \setlength{\gLinks}{\dimexpr\wd\gbox+6mm\relax}%
  \ifdim\gLinks>0.62\linewidth
    \usebox{\gbox}\par\smallskip\noindent\begin{minipage}{\linewidth}#2\end{minipage}%
  \else
    \begin{minipage}[t]{\gLinks}\vspace{0pt}\usebox{\gbox}\end{minipage}%
    \begin{minipage}[t]{\dimexpr\linewidth-\gLinks\relax}\vspace{0pt}\raggedright #2\end{minipage}%
  \fi\par}
% 16.10. (Lehrer, PA10 S. 1): Text links, Koordinatensystem rechts daneben; ist das System breiter als 0,58 Zeile,
% steht es wie bisher unter dem Text
% Passt das System knapp nicht mehr auf die Seite (bis 25 % zu hoch), wird es auf den Rest verkleinert statt
% die Seite halb leer zu lassen
\newlength{\gRest}\newlength{\gHoch}
\newcommand{\gtextks}[2]{\sbox{\gbox}{#1}%
  \setlength{\gHoch}{\dimexpr\ht\gbox+\dp\gbox\relax}%
  \ifdim\pagegoal<\maxdimen\setlength{\gRest}{\dimexpr\pagegoal-\pagetotal-4\baselineskip\relax}%
    \ifdim\gHoch>\gRest\ifdim\gRest>0.75\gHoch
      \sbox{\gbox}{\resizebox*{!}{\gRest}{\usebox{\gbox}}}\fi\fi\fi
  \ifdim\wd\gbox>0.58\linewidth #2\par\smallskip\noindent\gpassend{\usebox{\gbox}}\par
  \else\begin{minipage}[t]{\dimexpr\linewidth-\wd\gbox-5mm\relax}\raggedright #2\end{minipage}\hfill
    \begin{minipage}[t]{\wd\gbox}\vspace{0pt}\usebox{\gbox}\end{minipage}\par\fi}
% v0.9 (12.10., pruefung.md § 5 „Marke mit B, Bedingungen rechts“): drittes Argument von pfaufg = Bedingungen
% der Prüfung (v0.10: \pfohnetr = TR durchgestrichen, „FOR“), klein grau im rechten Rand auf der Grundlinie der Nummer, Zeilen mit \\
\newcommand{\gbedrand}[2]{\ifblank{#2}{}{\rlap{\hspace*{\dimexpr#1+2mm\relax}%
  \parbox[t]{15mm}{\raggedright\scriptsize\leavevmode\color{mbgrau}\strut#2}}}}
% Teil eines Bündels (a, b …): Marke links im Rand, Bedingungen rechts im Rand, auf der Zeile des Teils
\newcommand{\gteilrand}[2]{\llap{\mbpfrandmarke{#1}\hspace*{\mbpfnr}}\gbedrand{\linewidth}{#2}}
% Teil eines Bündels: Buchstabe in eigener Spalte, Inhalt eingerückt (\gEin = Einzug, damit die Spaltengrenze von
% \gzwei an derselben Stelle bleibt)
\newenvironment{gbteil}[3]{\par\smallskip\noindent\gteilrand{#2}{#3}\makebox[\mbpfteil][l]{\textbf{#1}}%
  \begin{minipage}[t]{\dimexpr\linewidth-\mbpfteil\relax}\setlength{\gEin}{\mbpfteil}\raggedright\strut\ignorespaces}%
 {\end{minipage}\par\addvspace{6pt}}% v0.18 (PG80 Nr. 1): Abstand, damit die Linie von a) nicht an b) rutscht
% Grafik höchstens so breit wie die Spalte (Bündel nebeneinander, Koordinatensystem ohne Karo)
\newsavebox{\gpbox}
\newcommand{\gpassend}[1]{\sbox{\gpbox}{#1}\ifdim\wd\gpbox>\linewidth\typeout{GPASSEND: verkleinert \the\wd\gpbox\space auf \the\linewidth}\resizebox{\linewidth}{!}{\usebox{\gpbox}}\else\usebox{\gpbox}\fi}
% Skizzenfeld (v0.10, Lehrer 13.10.): Karo ab dem linken Rand über die Zeile, „Skizze“ klein grau oben links im Feld
% (wie \pfskizzenfeld der Vorlage); der rechte Rand bleibt den Bedingungen
\newcommand{\gskizze}[1]{\par\smallskip\noindent\begin{tikzpicture}
  \clip (0,0) rectangle ({\linewidth-1mm},{#1*10mm});
  \draw[black!16,line width=0.3pt,step=5mm] (0,0) grid ({\linewidth},{#1*10mm});
  \node[anchor=north west,inner sep=2pt,font=\scriptsize,text=mbgrau,fill=white] at (0,{#1*10mm}) {Skizze};
  \end{tikzpicture}\par}
\RenewDocumentEnvironment{pfaufg}{m m m}{%
  \begin{lrbox}{\mbaufgabebox}\begin{minipage}[t]{\linewidth}%
  \gSchrift\noindent\hspace*{\gEin}\mbpfrandmarke{#1}\gbedrand{\linewidth-\gEin}{#3}%
  \begin{minipage}[t]{\mbpfnr}\strut\textbf{#2}\end{minipage}%
  \begin{minipage}[t]{\dimexpr\linewidth-\gEin-\mbpfnr\relax}\raggedright\strut\ignorespaces}%
 {\end{minipage}%
  \end{minipage}\end{lrbox}%
  \setlength{\mbaufgabehoehe}{\dimexpr\ht\mbaufgabebox+\dp\mbaufgabebox\relax}%
  \ifdefempty{\gkopfnach}{\par\addvspace{5pt}\Needspace*{\mbaufgabehoehe}}{%
    % Stufenkopf, der mit dieser (ersten) Aufgabe zusammenbleibt: Platz für beide verlangen, dann den Kopf setzen
    % (12.10., Quadratische: ein Kopf stand allein unten, die Aufgabe mit Koordinatensystem auf der nächsten Seite)
    % v0.17 (AS40 „Terme deuten“ allein unten): hinter dem Kopf steht das Sprungziel (\special, ein Knoten, nach
    % dem die Seite brechen durfte) – darum \nobreak nach dem Kopf; Vorhalt 56 pt = Kopf mit Linie und Abständen
    \par\Needspace*{\dimexpr\mbaufgabehoehe+56pt\relax}\gkopfnach\par\nobreak\gdef\gkopfnach{}\addvspace{5pt}}%
  \noindent\usebox{\mbaufgabebox}\par\addvspace{5pt}}
% v0.18 (plan.md § 10, 15.10. l): Abschnittsbuchstabe A, B, C vor dem Namen im Kopf – als Vorsatz, damit der
% Name im TeX gleich bleibt (Bestellkarte, Sprungziele lesen ihn)
\def\gbstakt{}\newcommand{\gbstvor}[1]{\gdef\gbstakt{#1}}
\NewCommandCopy\gpfstufekopfalt\pfstufekopf
\RenewDocumentCommand{\pfstufekopf}{o m m}{\IfValueTF{#1}{\gpfstufekopfalt[#1]}{\gpfstufekopfalt}%
  {\ifx\gbstakt\empty\else\gbstakt\enspace\fi #2}{#3}\gdef\gbstakt{}}
% v0.18 (16.10. e): Begriff grau, ohne Nummer, ohne Rahmen über der ersten Aufgabe der Stufe
\newcommand{\gbegriff}[1]{\par\nobreak\noindent\hspace*{\mbpfnr}\parbox[t]{\dimexpr\linewidth-\mbpfnr\relax}{\raggedright\small\color{mbgrau}#1}\par\nobreak\addvspace{3pt}}
% v0.18 (16.10. g): Lösungsblatt – Kopf mit kräftigem Strich, Abschnitt klein grau als Zwischenzeile
\newcommand{\glkopf}[1]{\noindent{\large\bfseries Lösungen \textperiodcentered{} #1}\par\vspace{3pt}\hrule height 1.2pt\par\medskip}
\newcommand{\glzwischen}[1]{\par\addvspace{8pt}\Needspace*{4\baselineskip}\noindent{\footnotesize\color{mbgrau}#1}\par\nobreak\vspace{1pt}}
% v0.18 (15.10. o, w): Blattende – Kern ohne Rahmen, dann graue Zeile „Mehr: n Aufgaben“
\newcommand{\gkern}[1]{\par\addvspace{14pt}\Needspace*{6\baselineskip}\noindent{\bfseries Kern}\par\nobreak\smallskip
  {\small\raggedright #1\par}}
\newcommand{\gmehr}[1]{\par\addvspace{10pt}\noindent{\footnotesize\color{mbgrau}Mehr: #1 Aufgaben}\par}
% Stufenkopf vormerken: \gkopfvor{<Kopf>} – die nächste pfaufg setzt ihn, sobald beide zusammen Platz haben
\def\gkopfnach{}
\newcommand{\gkopfvor}[1]{\gdef\gkopfnach{#1}}
% Teilaufgabe mit Skizze links in gzwei: Buchstabe in eigener Spalte, darunter Text und Skizze
\newcommand{\gteil}[2]{\makebox[\mbpfteil][l]{\textbf{#1}}\parbox[t]{\dimexpr0.5\textwidth-\gEin-\mbpfnr-5mm-\mbpfteil\relax}{\raggedright #2}}
% Stufenkopf der Vorlage (mathblatt.sty 2026-10-12: Linie, Abstand, 13 pt, \pfjahre) ohne ihren Vorhalt von
% 14 Zeilen: der Setzer fordert vor dem Kopf 0,16 Seitenhöhe (Kopf + erste Aufgabe), mehr nicht
\setlength{\mbpfstufevorhalt}{0pt}
\makeatother
\newcommand{\karo}[1]{\par\smallskip\noindent\begin{tikzpicture}
  \clip (0,0) rectangle ({\linewidth-1mm},{#1*10mm});
  \draw[black!16,line width=0.3pt,step=5mm] (0,0) grid ({\linewidth},{#1*10mm});
  \end{tikzpicture}\par}
\newcommand{\kk}[1]{$\square$\,#1\hspace{0.9em}}
\newcommand{\linie}[1]{\,{\color{black!45}\rule[-1pt]{#1}{0.4pt}}\,}
\newcommand{\pkt}{\textperiodcentered{}}
% v0.14 (14.10.): Stufenkopf mit fünf Punkten (Option = Jahre für die Bestellkarte) oder „zuletzt 19“
\NewDocumentCommand{\gpunkte}{o m}{#2}
\newcommand{\gzuletzt}[1]{zuletzt #1}
% v0.7 (12.10.): Schülerfrage unter dem Stufenkopf; Ankreuzoption als Bild (Buchstabe und Kästchen darunter)
\newcommand{\gfrage}[1]{\par\nobreak\noindent\hspace*{\mbpfnr}{\small\itshape #1}\par\nobreak\addvspace{2pt}}
\newlength{\gbildbreite}\setlength{\gbildbreite}{0.32\linewidth}% Setzer setzt sie je Reihe (vier Bilder: 0,24)
\newcommand{\gbildoption}[3]{\begin{minipage}[t][#3][b]{\gbildbreite}\centering #2\par\smallskip\mbox{#1\,$\square$}\end{minipage}}
% Lineare (12.10.): Bild mit Buchstabe ohne Kästchen (Diagramm zuordnen); Aussage mit „wahr / falsch“ zum Ankreuzen
\newcommand{\gbildname}[3]{\begin{minipage}[t][#3][b]{\gbildbreite}\centering #2\par\smallskip\mbox{\textbf{#1}}\end{minipage}}
\newcommand{\gwahrfalsch}[2]{\par\noindent\hangindent2.2em\hangafter1\makebox[2.2em][l]{(#1)}#2\hfill\mbox{wahr $\square$\quad falsch $\square$}\par}

\newcommand{\gunter}[2]{\par\addvspace{6pt}\noindent{\bfseries #1\enspace #2}\par\nobreak\vspace{2pt}}
\newcommand{\gsternzeile}{\par\vspace{-4pt}\noindent{\footnotesize\color{mbgrau}* nur FOR}\par\medskip}
% 17.10.: Skizze oben bündig mit der ersten Textzeile (beide Spalten hängen an der Oberkante); der Rechenraum (Karo)
% beginnt direkt unter dem Text, links neben dem Graph, und läuft unter dem Graph über die ganze Zeile weiter
% (Streifen links = auf 5 mm aufgerundete Höhe des Graphen unter dem Text, Raster durchgehend; die Fläche des Streifens
% wird von der Karohöhe unter dem Graph abgezogen, auf 0,5 cm aufgerundet – die Gesamtfläche Rechenraum bleibt mindestens #3 × Zeile); #3 = Zeilen Karo in cm
% unter dem Graph (0 = kein Karo). Kein Verkleinern des Graphen.
\newsavebox{\gtbox}\newlength{\gWl}\newlength{\gDh}
\newcommand{\gkaroohne}[1]{\nointerlineskip\noindent\begin{tikzpicture}
  \clip (0,0) rectangle ({\linewidth-1mm},{#1*10mm});
  \draw[black!16,line width=0.3pt,step=5mm] (0,0) grid ({\linewidth},{#1*10mm});
  \end{tikzpicture}\par}
\newcommand{\gtextkaro}[3]{\sbox{\gbox}{#1}%
  \ifdim\wd\gbox>0.58\linewidth #2\par\smallskip\noindent\gpassend{\usebox{\gbox}}\par\ifnum#3>0 \karo{#3}\fi
  \else
    \setlength{\gWl}{\dimexpr\linewidth-\wd\gbox-5mm\relax}%
    \sbox{\gtbox}{\begin{minipage}[t]{\gWl}\raggedright #2\par\end{minipage}}%
    \setlength{\gDh}{\dimexpr\ht\gbox+\dp\gbox-\ht\gtbox-\dp\gtbox-6pt\relax}%
    \leavevmode
    \ifnum#3>0\ifdim\gDh>0pt
      \pgfmathsetmacro{\gxs}{floor(\gWl/5mm)*5mm}\pgfmathsetmacro{\gys}{ceil(\gDh/5mm)*5mm}%
      \begin{minipage}[t]{\gWl}\usebox{\gtbox}\par\vspace{6pt}\nointerlineskip\noindent
        \begin{tikzpicture}
        \clip (0,0) rectangle (\gxs pt,\gys pt);
        \draw[black!16,line width=0.3pt,step=5mm] (0,0) grid (\gxs pt,\gys pt);
        \end{tikzpicture}\par\end{minipage}%
    \else\usebox{\gtbox}\fi\else\usebox{\gtbox}\fi
    \hfill\begin{minipage}[t]{\wd\gbox}\vspace{-\dimexpr\ht\strutbox+2.4ex\relax}\usebox{\gbox}\end{minipage}\par
    \ifnum#3>0 \ifdim\gDh>0pt
      \pgfmathsetmacro{\gnn}{max(0,ceil(2*(#3-(\gys/28.4528)*(\gxs/\linewidth)))/2)}%
      \ifdim\gnn pt>0pt\gkaroohne{\gnn}\fi
    \else\karo{#3}\fi\fi
  \fi}
\begin{document}

\pfheftstil
\pflaufkopf[P10 \pkt{} Prüfung]{Daten}{PS30 A}\pfseitevon
\pfheftkopf{Daten}{}
\pfstufenzeile{Kenngrößen}
\gsternzeile

% ===================== A Kenngrößen =====================
\gkopfvor{\gbstvor{A}\pfstufekopf{Kenngrößen}{}}
% DRUCK 2026-FOR-K3a deckt 2026-FOR-K3a 2026-EBR-K4a
% DRUCK 2026-FOR-K3b deckt 2026-FOR-K3b 2026-EBR-K4b
\begin{pfaufg}{nach~P10\\’26}{1.}{}
\gtextkaro{{\footnotesize\setlength{\tabcolsep}{4pt}\renewcommand{\arraystretch}{1.2}\begin{tabular}{|l|c|c|c|c|c|c|c|}\hline Wochentag & Mo & Di & Mi & Do & Fr & Sa & So\\\hline \begin{tabular}{@{}l@{}}Benzinpreis in\\€ pro Liter\end{tabular} & 1,77 & 1,78 & 1,77 & 1,79 & 1,84 & 1,82 & 1,85\\\hline\end{tabular}}}{An einer Tankstelle wurden die Benzinpreise eine Woche lang täglich um 18:00~Uhr notiert.\begin{gbteil}{a)}{}{}Gib die Spannweite der Benzinpreise an.\end{gbteil}}{3}
\par\medskip\noindent\hfill Spannweite $=$\linie{45mm}
\begin{gbteil}{b)}{}{}
Ein Kunde behauptet, dass die Tankstelle von Montag bis Freitag einen durchschnittlichen Benzinpreis von unter 1,80~€ pro Liter anbietet. Zeige rechnerisch, dass diese Aussage richtig ist.
\karo{4}
\end{gbteil}
\fusshilfe{\gfuss{1}{a) 0,08 € pro Liter; b) 1,79 € < 1,80 €}}
\end{pfaufg}
% DRUCK 2015-OS-K7b deckt 2015-OS-K7b
\begin{pfaufg}{nach~P10\\’15}{2.}{}
Eine Fußballmannschaft hat von insgesamt 34 Spielen 22-mal gewonnen, 10-mal unentschieden gespielt und 2-mal verloren. Gib die relative Häufigkeit an, mit der ein Spiel gewonnen wurde.
\karo{3}
\par\medskip\noindent\hfill relative Häufigkeit $=$\linie{45mm}
\fusshilfe{\gfuss{2}{0,65}}
\end{pfaufg}
% DRUCK 2025-OS-K6a deckt 2025-OS-K6a
\begin{pfaufg}{nach~P10\\’25}{3.}{}
Die Liste zeigt das Alter der Kinder einer Kita-Gruppe (in Jahren): 5 \quad 2 \quad 1 \quad 4 \quad 2 \quad 4 \quad 2 \quad 4 \quad 2 \quad 2. Kreuze die falsche Aussage an und korrigiere sie:\pfkreuzab{}{Die Spannweite beträgt 4 Jahre.}\pfkreuzab{}{Der Modalwert (Modus) beträgt 4 Jahre.}
\karo{3}
\par\medskip\noindent\hfill Modalwert $=$\linie{45mm}
\fusshilfe{\gfuss{3}{Modalwert falsch; 2 Jahre}}
\end{pfaufg}
% DRUCK 2024-OS-K2a deckt 2024-OS-K2a
\begin{pfaufg}{nach~P10\\’24}{4.}{}
\gtextkaro{\begin{tikzpicture}[font=\footnotesize,line width=0.6pt]\draw[black!14,line width=0.3pt] (0,0) -- (8,0);\draw[black!14,line width=0.3pt] (0,0.378) -- (8,0.378);\draw[black!14,line width=0.3pt] (0,0.756) -- (8,0.756);\draw[black!14,line width=0.3pt] (0,1.133) -- (8,1.133);\draw[black!14,line width=0.3pt] (0,1.511) -- (8,1.511);\draw[black!14,line width=0.3pt] (0,1.889) -- (8,1.889);\draw[black!14,line width=0.3pt] (0,2.267) -- (8,2.267);\draw[black!14,line width=0.3pt] (0,2.644) -- (8,2.644);\draw[black!14,line width=0.3pt] (0,3.022) -- (8,3.022);\draw[black!14,line width=0.3pt] (0,3.4) -- (8,3.4);\draw[black!40,line width=0.3pt] (0,0) -- (8,0);\node[left,inner sep=2pt] at (0,0) {0};\draw[black!40,line width=0.3pt] (0,0.378) -- (8,0.378);\node[left,inner sep=2pt] at (0,0.378) {10};\draw[black!40,line width=0.3pt] (0,0.756) -- (8,0.756);\node[left,inner sep=2pt] at (0,0.756) {20};\draw[black!40,line width=0.3pt] (0,1.133) -- (8,1.133);\node[left,inner sep=2pt] at (0,1.133) {30};\draw[black!40,line width=0.3pt] (0,1.511) -- (8,1.511);\node[left,inner sep=2pt] at (0,1.511) {40};\draw[black!40,line width=0.3pt] (0,1.889) -- (8,1.889);\node[left,inner sep=2pt] at (0,1.889) {50};\draw[black!40,line width=0.3pt] (0,2.267) -- (8,2.267);\node[left,inner sep=2pt] at (0,2.267) {60};\draw[black!40,line width=0.3pt] (0,2.644) -- (8,2.644);\node[left,inner sep=2pt] at (0,2.644) {70};\draw[black!40,line width=0.3pt] (0,3.022) -- (8,3.022);\node[left,inner sep=2pt] at (0,3.022) {80};\draw[black!40,line width=0.3pt] (0,3.4) -- (8,3.4);\node[left,inner sep=2pt] at (0,3.4) {90};\filldraw[fill=black!30,draw=black] (0.28,0) rectangle (0.68,1.972);\node[above,inner sep=1.5pt,font=\scriptsize\bfseries] at (0.48,1.972) {52,2};\node[rotate=45,anchor=north east,inner sep=2pt] at (0.63,0) {Januar};\filldraw[fill=black!30,draw=black] (0.92,0) rectangle (1.32,1.496);\node[above,inner sep=1.5pt,font=\scriptsize\bfseries] at (1.12,1.496) {39,6};\node[rotate=45,anchor=north east,inner sep=2pt] at (1.27,0) {Februar};\filldraw[fill=black!30,draw=black] (1.56,0) rectangle (1.96,1.334);\node[above,inner sep=1.5pt,font=\scriptsize\bfseries] at (1.76,1.334) {35,3};\node[rotate=45,anchor=north east,inner sep=2pt] at (1.91,0) {März};\filldraw[fill=black!30,draw=black] (2.2,0) rectangle (2.6,0.725);\node[above,inner sep=1.5pt,font=\scriptsize\bfseries] at (2.4,0.725) {19,2};\node[rotate=45,anchor=north east,inner sep=2pt] at (2.55,0) {April};\filldraw[fill=black!30,draw=black] (2.84,0) rectangle (3.24,1.862);\node[above,inner sep=1.5pt,font=\scriptsize\bfseries] at (3.04,1.862) {49,3};\node[rotate=45,anchor=north east,inner sep=2pt] at (3.19,0) {Mai};\filldraw[fill=black!30,draw=black] (3.48,0) rectangle (3.88,3.215);\node[above,inner sep=1.5pt,font=\scriptsize\bfseries] at (3.68,3.215) {85,1};\node[rotate=45,anchor=north east,inner sep=2pt] at (3.83,0) {Juni};\filldraw[fill=black!30,draw=black] (4.12,0) rectangle (4.52,1.364);\node[above,inner sep=1.5pt,font=\scriptsize\bfseries] at (4.32,1.364) {36,1};\node[rotate=45,anchor=north east,inner sep=2pt] at (4.47,0) {Juli};\filldraw[fill=black!30,draw=black] (4.76,0) rectangle (5.16,3.113);\node[above,inner sep=1.5pt,font=\scriptsize\bfseries] at (4.96,3.113) {82,4};\node[rotate=45,anchor=north east,inner sep=2pt] at (5.11,0) {August};\filldraw[fill=black!30,draw=black] (5.4,0) rectangle (5.8,0.82);\node[above,inner sep=1.5pt,font=\scriptsize\bfseries] at (5.6,0.82) {21,7};\node[rotate=45,anchor=north east,inner sep=2pt] at (5.75,0) {September};\filldraw[fill=black!30,draw=black] (6.04,0) rectangle (6.44,2.063);\node[above,inner sep=1.5pt,font=\scriptsize\bfseries] at (6.24,2.063) {54,6};\node[rotate=45,anchor=north east,inner sep=2pt] at (6.39,0) {Oktober};\filldraw[fill=black!30,draw=black] (6.68,0) rectangle (7.08,2.422);\node[above,inner sep=1.5pt,font=\scriptsize\bfseries] at (6.88,2.422) {64,1};\node[rotate=45,anchor=north east,inner sep=2pt] at (7.03,0) {November};\filldraw[fill=black!30,draw=black] (7.32,0) rectangle (7.72,1.59);\node[above,inner sep=1.5pt,font=\scriptsize\bfseries] at (7.52,1.59) {42,1};\node[rotate=45,anchor=north east,inner sep=2pt] at (7.67,0) {Dezember};\draw[-{Stealth[length=2mm]}] (0,0) -- (0,3.85) node[above right,inner sep=1pt] {Niederschlag in mm};\draw[-{Stealth[length=2mm]}] (0,0) -- (8.3,0);\node[above,font=\small\bfseries] at (4,4.15) {Niederschlag in Brandenburg 2021};\end{tikzpicture}}{Im Diagramm sind die monatlichen Niederschläge des Jahres 2021 in Brandenburg dargestellt. Ermittle die Spannweite der monatlichen Niederschläge in mm.}{3}
\par\medskip\noindent\hfill Spannweite $=$\linie{45mm}
\fusshilfe{\gfuss{4}{65,9 mm}}
\end{pfaufg}
% DRUCK 2020-OS-K2c deckt 2020-OS-K2c
% DRUCK 2020-OS-K2d deckt 2020-OS-K2d
\begin{pfaufg}{nach~P10\\’20}{5.}{}
\gtextkaro{\begin{tikzpicture}[font=\footnotesize,line width=0.6pt]\draw[black!14,line width=0.3pt] (0,0) -- (0,3.552);\draw[black!14,line width=0.3pt] (0.22,0) -- (0.22,3.552);\draw[black!14,line width=0.3pt] (0.44,0) -- (0.44,3.552);\draw[black!14,line width=0.3pt] (0.66,0) -- (0.66,3.552);\draw[black!14,line width=0.3pt] (0.88,0) -- (0.88,3.552);\draw[black!14,line width=0.3pt] (1.1,0) -- (1.1,3.552);\draw[black!14,line width=0.3pt] (1.32,0) -- (1.32,3.552);\draw[black!14,line width=0.3pt] (1.54,0) -- (1.54,3.552);\draw[black!14,line width=0.3pt] (1.76,0) -- (1.76,3.552);\draw[black!14,line width=0.3pt] (1.98,0) -- (1.98,3.552);\draw[black!14,line width=0.3pt] (2.2,0) -- (2.2,3.552);\draw[black!14,line width=0.3pt] (2.42,0) -- (2.42,3.552);\draw[black!14,line width=0.3pt] (2.64,0) -- (2.64,3.552);\draw[black!14,line width=0.3pt] (2.86,0) -- (2.86,3.552);\draw[black!14,line width=0.3pt] (3.08,0) -- (3.08,3.552);\draw[black!14,line width=0.3pt] (3.3,0) -- (3.3,3.552);\draw[black!14,line width=0.3pt] (3.52,0) -- (3.52,3.552);\draw[black!14,line width=0.3pt] (3.74,0) -- (3.74,3.552);\draw[black!14,line width=0.3pt] (3.96,0) -- (3.96,3.552);\draw[black!14,line width=0.3pt] (4.18,0) -- (4.18,3.552);\draw[black!14,line width=0.3pt] (4.4,0) -- (4.4,3.552);\draw[black!14,line width=0.3pt] (4.62,0) -- (4.62,3.552);\draw[black!14,line width=0.3pt] (4.84,0) -- (4.84,3.552);\draw[black!14,line width=0.3pt] (5.06,0) -- (5.06,3.552);\draw[black!14,line width=0.3pt] (5.28,0) -- (5.28,3.552);\draw[black!14,line width=0.3pt] (5.5,0) -- (5.5,3.552);\draw[black!14,line width=0.3pt] (5.72,0) -- (5.72,3.552);\draw[black!14,line width=0.3pt] (5.94,0) -- (5.94,3.552);\draw[black!14,line width=0.3pt] (6.16,0) -- (6.16,3.552);\draw[black!14,line width=0.3pt] (6.38,0) -- (6.38,3.552);\draw[black!14,line width=0.3pt] (6.6,0) -- (6.6,3.552);\draw[black!40,line width=0.3pt] (0,0) -- (0,3.552);\node[below,inner sep=2pt] at (0,0) {0};\draw[black!40,line width=0.3pt] (1.1,0) -- (1.1,3.552);\node[below,inner sep=2pt] at (1.1,0) {50};\draw[black!40,line width=0.3pt] (2.2,0) -- (2.2,3.552);\node[below,inner sep=2pt] at (2.2,0) {100};\draw[black!40,line width=0.3pt] (3.3,0) -- (3.3,3.552);\node[below,inner sep=2pt] at (3.3,0) {150};\draw[black!40,line width=0.3pt] (4.4,0) -- (4.4,3.552);\node[below,inner sep=2pt] at (4.4,0) {200};\draw[black!40,line width=0.3pt] (5.5,0) -- (5.5,3.552);\node[below,inner sep=2pt] at (5.5,0) {250};\draw[black!40,line width=0.3pt] (6.6,0) -- (6.6,3.552);\node[below,inner sep=2pt] at (6.6,0) {300};\filldraw[fill=black!30,draw=black] (0,0.196) rectangle (3.74,0.476);\node[left,inner sep=3pt] at (0,0.336) {Mo};\filldraw[fill=black!30,draw=black] (0,0.676) rectangle (2.86,0.956);\node[left,inner sep=3pt] at (0,0.816) {Di};\filldraw[fill=black!30,draw=black] (0,1.156) rectangle (4.62,1.436);\node[left,inner sep=3pt] at (0,1.296) {Mi};\node[left,inner sep=3pt] at (0,1.776) {Do};\filldraw[fill=black!30,draw=black] (0,2.116) rectangle (4.18,2.396);\node[left,inner sep=3pt] at (0,2.256) {Fr};\filldraw[fill=black!30,draw=black] (0,2.596) rectangle (5.28,2.876);\node[left,inner sep=3pt] at (0,2.736) {Sa};\filldraw[fill=black!30,draw=black] (0,3.076) rectangle (6.16,3.356);\node[left,inner sep=3pt] at (0,3.216) {So};\draw[-{Stealth[length=2mm]}] (0,0) -- (7,0);\draw (0,0) -- (0,3.552);\node[below,inner sep=2pt,yshift=-12pt] at (3.3,0) {Anzahl der Besucher};\node[rotate=90,above,inner sep=14pt] at (0,1.776) {Wochentag};\node[above,font=\small\bfseries] at (3.3,3.602) {Besucher im Schwimmbad};\end{tikzpicture}}{Im Diagramm ist die Anzahl der Besucher eines Schwimmbades in der ersten Ferienwoche der Sommerferien dargestellt. Am Donnerstag waren 180 Besucher im Schwimmbad.\begin{gbteil}{a)}{}{}Berechne die durchschnittliche Anzahl der Besucher in der ersten Ferienwoche.\end{gbteil}}{4}
\par\medskip\noindent\hfill Durchschnitt $=$\linie{45mm}
\begin{gbteil}{b)}{}{}
Die beiden folgenden Aussagen sind wahr. Weise sie rechnerisch nach. (1)~Die Spannweite der Besucherzahlen in der ersten Ferienwoche beträgt 150 Besucher. (2)~Am Sonntag kamen ein Drittel mehr Besucher als am Mittwoch.
\karo{5}
\end{gbteil}
\fusshilfe{\gfuss{5}{a) 200 Besucher; b) 150 Besucher; 280 Besucher}}
\end{pfaufg}
% DRUCK 2015-OS-K7a deckt 2015-OS-K7a
\begin{pfaufg}{nach~P10\\’15}{6.}{}
In der Saison 2012\slash{}13 hat Hertha BSC in der 2.~Bundesliga gespielt. Zu den 17 Heimspielen kamen insgesamt 680\,353 Zuschauer. Berechne, wie viele Zuschauer durchschnittlich ein Heimspiel besucht haben. Runde sinnvoll.
\karo{3}
\par\medskip\noindent\hfill Durchschnitt $=$\linie{45mm}
\fusshilfe{\gfuss{6}{40 000 Zuschauer}}
\end{pfaufg}
% DRUCK 2014-OS-K4a deckt 2014-OS-K4a
% DRUCK 2014-OS-K4b deckt 2014-OS-K4b
\begin{pfaufg}{nach~P10\\’14}{7.}{}
\gtextkaro{{\footnotesize\setlength{\tabcolsep}{4pt}\renewcommand{\arraystretch}{1.2}\begin{tabular}{|l|r|r|r|r|r|r|r|r|r|r|}\hline  & 03.07. & 10.07. & 17.07. & 24.07. & 31.07. & 07.08. & 14.08. & 21.08. & 28.08. & 04.09.\\\hline E10 & 154,2 & 157,4 & 158,8 & 159,7 & 159,1 & 161,0 & 164,8 & 169,2 & 168,1 & 167,0\\\hline Diesel & 139,6 & 142,5 & 144,2 & 145,4 & 146,2 & 147,1 & 150,7 & 154,0 & 153,1 & 152,0\\\hline\end{tabular}}}{In der Tabelle stehen die Preise an einer Tankstelle für die Kraftstoffe E10 und Diesel in Cent pro Liter, vom 03.07. bis zum 04.09.2012.}{0}
\begin{gbteil}{a)}{}{}
Gib das Minimum und das Maximum des Preises für E10 in diesem Zeitraum an.
\par\medskip\noindent\hfill Minimum $=$\linie{25mm}\quad Maximum $=$\linie{25mm}
\end{gbteil}
\begin{gbteil}{b)}{}{}
Berechne die Spannweite des Preises für Diesel in diesem Zeitraum.
\karo{3}
\par\medskip\noindent\hfill Spannweite $=$\linie{45mm}
\end{gbteil}
\fusshilfe{\gfuss{7}{a) 154,2 ct; 169,2 ct; b) 14,4 ct}}
\end{pfaufg}
% DRUCK 2023-OS-K6c deckt 2023-OS-K6c
\begin{pfaufg}{nach~P10\\’23}{8.}{}
\gtextkaro{{\footnotesize\setlength{\tabcolsep}{4pt}\renewcommand{\arraystretch}{1.2}\begin{tabular}{|l|r|}\hline \textbf{Sprache} & \begin{tabular}{@{}c@{}}\textbf{Anzahl der Menschen}\\\textbf{mit dieser Muttersprache}\\\textbf{(in Millionen)}\end{tabular}\\\hline Arabisch & 290\\\hline Chinesisch & 1\,300\\\hline Englisch & 500\\\hline Hindi & 525\\\hline Spanisch & 389\\\hline\end{tabular}}}{Es gibt vermutlich mehr als 7\,000 Sprachen auf der Welt. In der Tabelle sind die fünf meistgesprochenen Sprachen angegeben. Gib Minimum, Maximum und Spannweite der Datenreihe an.}{3}
\par\medskip\noindent\hfill Minimum $=$\linie{25mm}\quad Maximum $=$\linie{25mm}\quad Spannweite $=$\linie{25mm}
\fusshilfe{\gfuss{8}{290 Mio.; 1300 Mio.; 1010 Mio.}}
\end{pfaufg}
% DRUCK 2022-OS-K4a deckt 2022-OS-K4a
\begin{pfaufg}{nach~P10\\’22}{9.}{}
\gtextkaro{\begin{tikzpicture}[font=\footnotesize,line width=0.6pt]\draw[black!14,line width=0.3pt] (0,0) -- (6.325,0);\draw[black!14,line width=0.3pt] (0,0.364) -- (6.325,0.364);\draw[black!14,line width=0.3pt] (0,0.727) -- (6.325,0.727);\draw[black!14,line width=0.3pt] (0,1.091) -- (6.325,1.091);\draw[black!14,line width=0.3pt] (0,1.455) -- (6.325,1.455);\draw[black!14,line width=0.3pt] (0,1.818) -- (6.325,1.818);\draw[black!14,line width=0.3pt] (0,2.182) -- (6.325,2.182);\draw[black!14,line width=0.3pt] (0,2.545) -- (6.325,2.545);\draw[black!14,line width=0.3pt] (0,2.909) -- (6.325,2.909);\draw[black!14,line width=0.3pt] (0,3.273) -- (6.325,3.273);\draw[black!14,line width=0.3pt] (0,3.636) -- (6.325,3.636);\draw[black!14,line width=0.3pt] (0,4) -- (6.325,4);\draw[black!40,line width=0.3pt] (0,0) -- (6.325,0);\node[left,inner sep=2pt] at (0,0) {390};\draw[black!40,line width=0.3pt] (0,0.364) -- (6.325,0.364);\node[left,inner sep=2pt] at (0,0.364) {400};\draw[black!40,line width=0.3pt] (0,0.727) -- (6.325,0.727);\node[left,inner sep=2pt] at (0,0.727) {410};\draw[black!40,line width=0.3pt] (0,1.091) -- (6.325,1.091);\node[left,inner sep=2pt] at (0,1.091) {420};\draw[black!40,line width=0.3pt] (0,1.455) -- (6.325,1.455);\node[left,inner sep=2pt] at (0,1.455) {430};\draw[black!40,line width=0.3pt] (0,1.818) -- (6.325,1.818);\node[left,inner sep=2pt] at (0,1.818) {440};\draw[black!40,line width=0.3pt] (0,2.182) -- (6.325,2.182);\node[left,inner sep=2pt] at (0,2.182) {450};\draw[black!40,line width=0.3pt] (0,2.545) -- (6.325,2.545);\node[left,inner sep=2pt] at (0,2.545) {460};\draw[black!40,line width=0.3pt] (0,2.909) -- (6.325,2.909);\node[left,inner sep=2pt] at (0,2.909) {470};\draw[black!40,line width=0.3pt] (0,3.273) -- (6.325,3.273);\node[left,inner sep=2pt] at (0,3.273) {480};\draw[black!40,line width=0.3pt] (0,3.636) -- (6.325,3.636);\node[left,inner sep=2pt] at (0,3.636) {490};\draw[black!40,line width=0.3pt] (0,4) -- (6.325,4);\node[left,inner sep=2pt] at (0,4) {500};\filldraw[fill=black!30,draw=black] (0.562,0) rectangle (1.162,2.836);\node[above,inner sep=1.5pt,font=\scriptsize\bfseries] at (0.862,2.836) {468};\node[below,inner sep=3pt] at (0.862,0) {2017};\filldraw[fill=black!30,draw=black] (1.712,0) rectangle (2.312,1.527);\node[above,inner sep=1.5pt,font=\scriptsize\bfseries] at (2.012,1.527) {432};\node[below,inner sep=3pt] at (2.012,0) {2018};\filldraw[fill=black!30,draw=black] (2.862,0) rectangle (3.462,3.455);\node[above,inner sep=1.5pt,font=\scriptsize\bfseries] at (3.162,3.455) {485};\node[below,inner sep=3pt] at (3.162,0) {2019};\filldraw[fill=black!30,draw=black] (4.013,0) rectangle (4.612,2.909);\node[above,inner sep=1.5pt,font=\scriptsize\bfseries] at (4.312,2.909) {470};\node[below,inner sep=3pt] at (4.312,0) {2020};\filldraw[fill=black!30,draw=black] (5.162,0) rectangle (5.762,0.364);\node[above,inner sep=1.5pt,font=\scriptsize\bfseries] at (5.462,0.364) {400};\node[below,inner sep=3pt] at (5.462,0) {2021};\draw[-{Stealth[length=2mm]}] (0,0) -- (0,4.45) node[above right,inner sep=1pt] {Anzahl der Jugendlichen};\draw[-{Stealth[length=2mm]}] (0,0) -- (6.625,0) node[below left,inner sep=2pt,yshift=-9pt] {Jahr};\node[above,font=\small\bfseries] at (3.162,4.75) {Bücher lesen von 2017 bis 2021};\end{tikzpicture}}{In den Jahren 2017 bis 2021 wurden jeweils 1\,200 Jugendliche in Deutschland befragt, wie oft sie Bücher lesen. Das Diagramm zeigt, wie viele der Befragten mehrmals pro Woche Bücher lesen. Gib das Minimum und das Maximum der Anzahl der Jugendlichen an, die mehrmals pro Woche Bücher lesen. Bestimme die durchschnittliche Anzahl (das arithmetische Mittel) dieser Jugendlichen.}{4}
\par\medskip\noindent\hfill Minimum $=$\linie{25mm}\quad Maximum $=$\linie{25mm}\quad Mittelwert $=$\linie{25mm}
\fusshilfe{\gfuss{9}{400; 485; 451}}
\end{pfaufg}
% DRUCK 2025-OS-K6c deckt 2025-OS-K6c
\begin{pfaufg}{nach~P10\\’25}{10.\textsuperscript{*}}{}
Die Liste zeigt das Alter der 11 Erziehenden in einer Kita: 52 \quad 57 \quad 33 \quad 44 \quad 31 \quad 30 \quad 66 \quad 58 \quad 61 \quad 41 \quad 55. Bestimme das Durchschnittsalter der Erziehenden. Die älteste und die jüngste Person verlassen die Kita. Erkläre, wie sich diese Veränderung auf das Durchschnittsalter auswirkt.
\karo{6}
\par\medskip\noindent\hfill Durchschnittsalter $=$\linie{45mm}
\fusshilfe{\gfuss{10}{48 Jahre; bleibt gleich}}
\end{pfaufg}
\end{document}